\begin{sloppypar}
Odor plumes in nature are turbulent and intermittent, forcing navigating animals to continually arbitrate between exploratory search (to re-locate the plume) and stable goal-directed locomotion (to navigate up the plume toward the odor source). Here we combine a closed-loop virtual plume paradigm in Drosophila with circuit mapping, population imaging and optogenetic perturbations to identify a neural mechanism for this arbitration. Flies tracking a plume exhibit two separable locomotor modes: sparse odor encounters evoke slow exploratory circling, whereas frequent odor encounters promote faster, straighter runs aligned to the long axis of the plume. Activity in two mushroom body output neurons (MBONs) forms an opponent representation of recent odor encounter statistics, with an increase in odor encounter frequency producing higher MBON21 activity and lower MBON09 activity. These MBONs then modulate the brain’s internal representation of travel direction in hΔB cells of the central complex. Activating MBON21 suppresses turning and increases path straightness, whereas activating hΔB biases turning away from the fly’s pre-stimulus goal heading. Together, these results identify an MBON--central complex pathway that converts odor encounter history into navigation control, enabling robust navigation in intermittent sensory environments.
\end{sloppypar}
